% Original Questions and Inferences, printed 30--43; previous packet frozen.
\section{Erotetic search scenarios and search rules}
\label{sec:p18-erotetic-search-scenarios}

A warranted auxiliary question is not yet an inquiry plan. A plan must say
which question is actually asked, what happens on every direct response, and
how a final answer follows from the information collected on that branch.
Wi\'sniewski's search scenarios supply this additional structure. Their native
objects are finite families of interconnected erotetic derivations, representable
as trees \cite[secs.~9--10]{Wisniewski2001QuestionsInferences}. A question can
occur without being asked: it can justify a later query. Removing such a
question may produce useful instructions while destroying the derivation.

Use the fixed declarative language, admissible partitions, direct-answer
assignment and semantic relations of
Definitions~\ref{def:p18-ei-setting} and~\ref{def:p18-ei-relation}.
Soundness means that at least one direct answer is true, not that a question
itself has a truth value. The current section is not a new PMCE calculus.
\begin{hbconvention}{def:p18-es-native-notation}{Native classical and question notation}
The source writes \(\supset\) for material implication and \(\equiv\)
for the classical biconditional. Here these are object-language connectives,
not entailment or elementary equivalence. Its connective-schema parameter
\(\otimes\) ranges over \(\land,\lor,\supset,\equiv\); it is not the
linear-logic tensor. The marker \(\pm\) in \(?\pm|A,B|\) identifies the
binary conjunctive question whose four direct answers are displayed below.
It is not a signed numerical error bound.
\end{hbconvention}

In the Boolean examples, abbreviate \(?\{A,\neg A\}\) as \(?A\), and write
\[
 ?\pm|A,B|=?\{A\land B,A\land\neg B,
                         \neg A\land B,\neg A\land\neg B\}.
\]
This is the source's binary conjunctive question. All displayed question
expressions must be well formed, with syntactically distinct direct answers.

\subsection{Derivations, queries and fair path families}

\begin{hbdefinition}{def:p18-es-derivation}{A source-native erotetic derivation}
Let \(X\) be a set of declaratives and \(Q\) a question. A finite sequence
\(e=(\varphi_1,\ldots,\varphi_n)\) is an erotetic derivation of
\(A\in dQ\) on \(X\) when \(\varphi_1=Q\), \(\varphi_n=A\), and:
\begin{enumerate}
\item Every question \(\varphi_k\) with \(k>1\) has
\(d\varphi_k\ne dQ\), and its next element is either a question or
a direct answer to \(\varphi_k\).
\item Every declarative \(\varphi_j\) belongs to \(X\), or is a direct
answer to the immediately previous question, which is not \(Q\), or is
entailed by a set of declaratives all occurring earlier in the sequence.
\item Every question other than \(Q\) is erotetically implied by some
earlier question, on the empty basis or on a set of declaratives all occurring
earlier in the sequence.
\end{enumerate}
Earlier does not mean immediately earlier. A \emph{query occurrence} is
an index \(k\) with \(1<k<n\) where \(\varphi_k\) is a question and
\(\varphi_{k+1}\in d\varphi_k\). An auxiliary question followed by
another question is not queried there: it serves as an erotetic premise.
These are Definition 12 and its query convention in the original source
\cite[pp.~34--35]{Wisniewski2001QuestionsInferences}.
\end{hbdefinition}

\begin{hbdefinition}{def:p18-es-family}{The source's search-scenario family conditions}
A finite family \(\Phi\) of at least two such derivations is a search
scenario for \(Q\) relative to \(X\) when \(dQ\cap X=\varnothing\)
and the following two conditions hold at every proper prefix of every path:
\begin{enumerate}
\item If the current element is a question and the next element is one of
its direct answers, every direct answer to that question is the next element
of some path sharing the entire prefix through the question.
\item If the current element is declarative, or is a question whose next
element is not a direct answer, every path sharing that prefix has the same
next element.
\end{enumerate}
These successor comparisons require that the indicated successor exists.
Thus paths are terminal paths, not proper prefixes of other paths. The
first condition is response fairness, not a probability or truth assumption.
It includes the index 1 in the literal source definition, although the query
convention above excludes that index. A scenario is \emph{complete} when
every member of \(dQ\) is an endpoint of some path; otherwise it is incomplete
\cite[p.~35]{Wisniewski2001QuestionsInferences}.
\end{hbdefinition}

\begin{hbdefinition}{def:p18-es-tree}{Prefix tree and operational root separation}
The tree of \(\Phi\) has one node for each occurring finite prefix; an
edge appends one expression. The root is \((Q)\), a node is labelled by its
last expression, and terminal nodes are the complete paths. The same formula
label can occur at several nodes; those are different inquiry states.
A branching occurrence has two different immediate successors. At indices
greater than one this is the source's branching-point definition.

For an operational scenario impose \emph{root separation}: no path has
a direct answer to \(Q\) as its second element. This extra convention is
stated explicitly rather than silently incorporated into Definition 13.
The source's ordinary Figures 1--3 satisfy it.
\end{hbdefinition}

\begin{hbtheorem}{thm:p18-es-structure}{proved=annex:p18.es.structure}{Tree structure and finite response branching}
Only questions are queries. Every query is a branching occurrence, and no
declarative occurrence branches. Each queried answer set is finite, with
at most \(|\Phi|\) distinct answers. In a root-separated scenario there is
a common first query on every path; it is the first branching occurrence.
The tree is finite, with height bounded by the largest path length.
An infinite-answer question can occur as an unqueried premise, but cannot
be a fair query of this finite scenario definition.
\end{hbtheorem}

\begin{hbtheorem}{thm:p18-es-definition-boundaries}{proved=annex:p18.es.definition-boundaries}{Two literal-definition edge cases}
If empty supporting sets are permitted in the derivability clause, the literal
definitions admit a two-path family \((Q,A),(Q,B)\), where
\(dQ=\{A,B\}\) and \(A,B\) are syntactically distinct Boolean tautologies.
It has no query under the source's \(k>1\) convention. Hence the source's
unqualified Facts 6--7 need root separation or a further explicit restriction.

The source's move-only canonical recipe also needs care when a query answer
belongs to \(X\). Such an answer occurrence cannot simply be moved away
from its query while preserving the scenario. A complete counterfamily is
supplied in the annex; this does not claim that no different canonical
scenario could ever be constructed.
\end{hbtheorem}

\subsection{The Golden Path guarantee}

\begin{hbtheorem}{thm:p18-es-golden-path}{proved=annex:p18.es.golden-path}{The generalized Golden Path Theorem}
Let \(\Phi\) be a source-defined scenario for \(Q\) relative to \(X\).
Fix an admissible partition at which every member of \(X\) is true and
\(Q\) is sound. Then \(\Phi\) contains a path on which every declarative
is true, every question is sound, and the final direct answer to \(Q\) is true.
The scenario need not be complete. The proof also handles the literal
root-answer edge case rather than relying on root separation.
\end{hbtheorem}

The original states this theorem and refers to the propositional proof in
a separate 2003 paper, saying that the generalized proof proceeds similarly
\cite[p.~36]{Wisniewski2001QuestionsInferences}. The annex supplies a complete
generalized proof from the actual family conditions; it does not treat that
unread referral as an imported completed derivation.

\begin{hbcorollary}{thm:p18-es-execution}{proved=annex:p18.es.execution}{Conditional finite execution, not guaranteed truthful access}
If a root-separated finite scenario is executed on this fixed partition,
every requested query returns a true direct answer, and the prescribed
declarative consequences are correctly obtained, then each resulting run
ends at a true direct answer to \(Q\). It uses at most the tree height minus
one successor steps. The condition does not guarantee that an oracle will
respond, that an accepted response is true, that consequence checking is
cheap, or that a finite scenario can be found for every inquiry.
\end{hbcorollary}

Soundness is existential: a sound query has a true answer in the model, not
necessarily an accessible answer. Response fairness records all answers,
including false and mutually inconsistent responses on different paths.
There may be more than one true response; the guarantee permits any truthful
one. A finite branch diagram is therefore not a proof of practical time
termination when a query can remain unanswered.

\subsection{Canonical, concise and imperative forms}

\begin{hbdefinition}{def:p18-es-transforms}{The three source transformations}
A scenario is canonical when no formula belonging to \(X\) occurs after
the first branching occurrence. It is concise when it is canonical and
each declarative occurrence is an initial premise, an immediately following
query answer, or a final direct answer to \(Q\).
The imperative counterpart of a concise scenario deletes its auxiliary
question occurrences that are not queries. It retains the principal question,
premises, queried questions, responses and final answers
\cite[pp.~36--37]{Wisniewski2001QuestionsInferences}.
\end{hbdefinition}

\begin{hblemma}{thm:p18-es-basis-replacement}{proved=annex:p18.es.basis-replacement}{Replacing derived auxiliary premises by their retained grounds}
If \(Z\models_L D\) for every \(D\in Y\), then
\(\operatorname{Im}_L(Q,Y,R)\) implies
\(\operatorname{Im}_L(Q,Z,R)\). Each original proper answer-set witness
is retained. This is the semantic fact needed when intermediate declarative
consequences are removed; it is not unrestricted transitivity of questions.
\end{hblemma}

\begin{hbtheorem}{thm:p18-es-canonical-concise}{proved=annex:p18.es.canonical-concise}{Explicit canonical and concise constructions}
Let the scenario be root-separated and suppose the direct-answer set of
every query is disjoint from \(X\). Collect the finitely many initial
premises actually occurring on its paths. Place them, in one fixed order,
immediately after the root on every path, and remove their later occurrences.
The result is a canonical scenario. This condition prevents removing a
required response edge merely because its formula belongs to \(X\).

For any canonical operational scenario, delete the intermediate declaratives that are
neither initial premises, immediate query responses nor final answers.
The result is concise and remains a scenario. The proof replaces every
deleted premise by its retained grounds and checks response fairness,
deterministic prefixes and endpoint preservation.
\end{hbtheorem}

\begin{hbtheorem}{thm:p18-es-imperative}{proved=annex:p18.es.imperative}{What survives imperative projection}
The imperative counterpart of a concise operational scenario retains a
truthful path to a true direct answer under the Golden Path assumptions.
It need not itself be a scenario: a retained query may lose the earlier
question that supplied its implication warrant. The predicate example below
and the connective procedures give explicit failures of that warrant after
projection. Correct conditional instructions and a valid erotetic derivation
are distinct objects.
\end{hbtheorem}

\subsection{Original scenario: meeting and talking}

\begin{hbexample}{def:p18-es-meeting}{Three direct answers, three response paths}
Let \(p\) mean John met Mary and \(q\) that he talked to her. The principal
question is construed with \(dQ=\{\neg p,p\land q,p\land\neg q\}\),
not merely as a polar question about a conjunction. The scenario is
\[
\begin{array}{l}
 (Q,?p,\neg p),\\
 (Q,?p,p,?q,q,p\land q),\\
 (Q,?p,p,?q,\neg q,p\land\neg q).
\end{array}
\]
Ask whether the meeting occurred. A negative answer is already a requested
answer. On a positive answer ask whether talking occurred, then combine
the two responses. Every response is represented, and every principal
direct answer is an endpoint: this scenario is complete.

The principal answers and terminal branches agree with the source's
prose and schematic meeting procedure
\cite[pp.~31,38]{Wisniewski2001QuestionsInferences}.
\end{hbexample}

\begin{hbtheorem}{thm:p18-es-meeting}{proved=annex:p18.es.meeting}{Every warrant and response in the meeting scenario}
The displayed family is a complete, canonical, concise scenario on the empty
basis. Its first query follows from \(Q\); on basis \(\{p\}\), its second
query follows from the earlier \(Q\). Neither pure implication between
independent \(?p\) and \(?q\) nor unrestricted question transitivity is used.
Every response branch has the stated direct endpoint.
\end{hbtheorem}

\subsection{Original scenario: departure and umbrella}

\begin{hbexample}{def:p18-es-destination}{A fair but incomplete scenario}
Use \(p,q,r,s\) for Paris, London, Kiev, Moscow, \(t\) for a morning
departure, and \(u\) for taking the umbrella. Put
\[
 Q=?\{p,q,r,s\},\quad
 X=\{x_1=t\supset(p\lor q),\ x_2=\neg t\supset s,\ x_3=q\equiv u\}.
\]
Here \(x_i=\cdots\) names the formulas; equality is not an extra premise.
With \(R=?\{p,q\}\), the three source paths are
\[
\begin{array}{l}
 (Q,x_1,x_2,x_3,?t,\neg t,s),\\
 (Q,x_1,x_2,x_3,?t,t,R,?u,u,q),\\
 (Q,x_1,x_2,x_3,?t,t,R,?u,\neg u,p).
\end{array}
\]
The first query distinguishes morning from nonmorning. On the nonmorning
answer, Moscow follows. On morning, the unqueried \(R\) warrants asking
about the umbrella. Its positive response yields London; its negative response
and morning coverage yield Paris. The Kiev answer is not an endpoint, so this
is incomplete under the source's definition, although both responses to each
query are represented \cite[pp.~31--32,35]{Wisniewski2001QuestionsInferences}.
\end{hbexample}

\begin{hbtheorem}{thm:p18-es-destination}{proved=annex:p18.es.destination}{All departure branches and the omitted endpoint}
The displayed family is a canonical, concise scenario relative to \(X\).
Its three endpoints follow on their respective branches; all question warrants
and response-fairness obligations hold. The absence of endpoint \(r\) does
not invalidate Golden Path. No exclusive-destination axiom or claim that
\(r\) is false is silently added to the propositional reconstruction.
\end{hbtheorem}

\subsection{Original predicate scenario and its imperative counterpart}

\begin{hblemma}{thm:p18-es-binary-warrants}{proved=annex:p18.es.binary-warrants}{The binary question hidden behind the instructions}
For each Boolean connective \(\otimes\in\{\land,\lor,\supset,\equiv\}\),
\(?(A\otimes B)\) regularly implies \(?\pm|A,B|\) on the empty
basis; \(?\pm|A,B|\) implies \(?A\) and \(?B\).
For independent variables \(A,B\), the direct pure implication from
\(?(A\otimes B)\) to \(?A\) fails in each of the four cases.
These are the actual two-leg warrants in note 15, not a transitivity shortcut
\cite[pp.~39--40]{Wisniewski2001QuestionsInferences}.
\end{hblemma}

\begin{hbexample}{def:p18-es-predicate}{An unasked binary question is still a proof premise}
Let \(P_1,P_2,P_3\) be distinct predicate symbols. Let \(Q=?P_1a\),
let \(H=\forall x(P_1x\equiv(P_2x\land P_3x))\),
and put \(R=?\pm|P_2a,P_3a|\). The complete source scenario has paths
\[
\begin{array}{l}
 (Q,H,R,?P_2a,\neg P_2a,\neg P_1a),\\
 (Q,H,R,?P_2a,P_2a,?P_3a,P_3a,P_1a),\\
 (Q,H,R,?P_2a,P_2a,?P_3a,\neg P_3a,\neg P_1a).
\end{array}
\]
The universal premise maps each binary conjunctive answer to the appropriate
principal polar answer. The unqueried \(R\) warrants both actual queries.
A negative \(P_2a\) stops the inquiry with a negative principal answer;
a positive response makes \(P_3a\) the next query. Both responses to that
query give the shown endpoints \cite[pp.~33--34]{Wisniewski2001QuestionsInferences}.

The imperative counterpart deletes \(R\) from these three paths, exactly
as Figure 4 does. It remains a correct conditional instruction set, but the
positive \(P_2a\) response alone does not reduce the principal polar alternatives:
\(P_3a\) can still be true or false. Thus the projected first query has lost
its implication warrant. The actual earlier binary question supplies that
warrant; pure implication between the independent polar queries does not.
\end{hbexample}

\begin{hbtheorem}{thm:p18-es-predicate}{proved=annex:p18.es.predicate}{Predicate correctness and failure of the projected derivation}
The full family is a complete, canonical, concise scenario relative to
\(\{H\}\), and every listed response branch is correct. Its imperative
counterpart satisfies the projected Golden Path conclusion but is not an
erotetic derivation family under Definition~\ref{def:p18-es-derivation}.
One-element first-order models witness the missing implication warrant.
\end{hbtheorem}

\subsection{The source's search-rule procedures, with every response}

The first three schematic procedures, Figures 5--7, abstract the preceding
scenarios: the meeting schema has direct answers
\(\{\neg A,A\land B,A\land\neg B\}\); the departure schema has basis
\(\{E\supset(A\lor B),\neg E\supset D,B\equiv F\}\) and principal
answers \(\{A,B,C,D\}\); the predicate schema uses
\(\forall x(Px\equiv(Rx\land Tx))\) for principal \(?Pc\).
Their exact instruction and response branches appear in the sheet. Different
source variables/nonlogical symbols are assigned distinct metavariables;
well-formedness is checked for each instance.

Figure 8 asks \(?A\) in order to answer \(?\neg A\). Response \(A\)
yields \(\neg\neg A\); response \(\neg A\) is already a principal answer.
Figures 9--12 instead answer a polar compound question. For conjunction,
ask \(?A\), stop negatively on \(\neg A\), and otherwise ask \(?B\).
For implication, stop positively on \(\neg A\), and otherwise ask \(?B\).
For disjunction, stop positively on \(A\), and otherwise ask \(?B\).
For equivalence ask \(?B\) on either first response: agreement yields
\(A\equiv B\), disagreement yields its negation.

In each compound case the original underlying scenario inserts the unasked
\(?\pm|A,B|\) between the principal question and \(?A\). It also
supplies the warrant for any later \(?B\). The imperative diagrams omit
that premise, as note 15 explicitly says; the sheet is an instruction sheet,
not a new set of direct question-implication axioms.

\calculussheet{sheet:p18-erotetic-search}{Original erotetic search warrants and imperative procedures}
  {E/Im entries are source warrant templates, not primitive axioms. The eight
  procedure entries reproduce every response case of Figures 5--12, not direct
  question-implication rules. Full compound scenarios retain the unqueried
  \(?\pm|A,B|\). Instances preserve question well-formedness.
  Proofs: \ref{annex:p18.es.search-procedures}.}

\hbrule{rule:p18-es-e-warrant}{E-rule warrant (source schema)}
  {X\qquad E_L(X,Q)}{\text{raise }Q}
  {\(X\) finite; actual evocation gives coverage and individual non-entailment.
   Actual soundness requires true \(X\).}

\hbrule{rule:p18-es-im-warrant}{Im-rule warrant (source schema)}
  {Q\qquad X\qquad\operatorname{Im}_L(Q,X,R)}{\text{raise }R}
  {Actual two-clause relation with proper nonempty witnesses; requires sound
   Q and true X for soundness. Earlier question need not have been queried.}

\hbrule{rule:p18-es-meeting}{Figure 5: meeting (imperative)}
  {Q=?\{\neg A,A\land B,A\land\neg B\}}
  {\text{ask }?A;\quad\begin{array}{l}\neg A\mapsto\neg A\\
   A\mapsto\text{ask }?B;\quad B\mapsto A\land B,\ \neg B\mapsto A\land\neg B\end{array}}
  {Empty basis. Second query uses original Q and accepted A.
   All three branches: \ref{annex:p18.es.meeting}.}

\hbrule{rule:p18-es-departure}{Figure 6: departure (imperative)}
  {Q=?\{A,B,C,D\};\quad E\supset(A\lor B),\ \neg E\supset D,\ B\equiv F}
  {\text{ask }?E;\quad\begin{array}{l}\neg E\mapsto D\\
   E\mapsto\text{ask }?F;\quad F\mapsto B,\ \neg F\mapsto A\end{array}}
  {Full positive-E branch has unqueried \(?\{A,B\}\).
   No C-endpoint: incomplete but fair. \ref{annex:p18.es.destination}.}

\hbrule{rule:p18-es-predicate}{Figure 7: predicate (imperative)}
  {?Pc;\quad\forall x(Px\equiv(Rx\land Tx))}
  {\text{ask }?Rc;\quad\begin{array}{l}\neg Rc\mapsto\neg Pc\\
   Rc\mapsto\text{ask }?Tc;\quad Tc\mapsto Pc,\ \neg Tc\mapsto\neg Pc\end{array}}
  {Distinct predicates; closed instances; nonempty domain. Full scenario has
   \(?\pm|Rc,Tc|\). \ref{annex:p18.es.predicate}.}

\hbrule{rule:p18-es-negation}{Figure 8: negation (imperative)}
  {?\neg A}
  {\text{ask }?A;\quad A\mapsto\neg\neg A,\ \neg A\mapsto\neg A}
  {Boolean double negation. Both responses: \ref{annex:p18.es.search-procedures}.}

\hbrule{rule:p18-es-conjunction}{Figure 9: conjunction (imperative)}
  {?(A\land B)}
  {\text{ask }?A;\quad\begin{array}{l}\neg A\mapsto\neg(A\land B)\\
   A\mapsto\text{ask }?B;\quad B\mapsto A\land B,\ \neg B\mapsto\neg(A\land B)\end{array}}
  {Boolean truth; full scenario retains binary question; no pure shortcut to
   \(?A\). \ref{annex:p18.es.search-procedures}.}

\hbrule{rule:p18-es-implication}{Figure 10: implication (imperative)}
  {?(A\supset B)}
  {\text{ask }?A;\quad\begin{array}{l}\neg A\mapsto A\supset B\\
   A\mapsto\text{ask }?B;\quad B\mapsto A\supset B,\ \neg B\mapsto\neg(A\supset B)\end{array}}
  {Boolean material conditional. Full scenario retains binary question.
   Negative-A branch is positive. \ref{annex:p18.es.search-procedures}.}

\hbrule{rule:p18-es-disjunction}{Figure 11: disjunction (imperative)}
  {?(A\lor B)}
  {\text{ask }?A;\quad\begin{array}{l}A\mapsto A\lor B\\
   \neg A\mapsto\text{ask }?B;\quad B\mapsto A\lor B,\ \neg B\mapsto\neg(A\lor B)\end{array}}
  {Boolean truth; full scenario retains binary question. All three routes:
   \ref{annex:p18.es.search-procedures}.}

\hbrule{rule:p18-es-equivalence}{Figure 12: equivalence (imperative)}
  {?(A\equiv B)}
  {\text{ask }?A;\ \text{then ask }?B;\quad\begin{array}{ll}
   A,B\mapsto A\equiv B&A,\neg B\mapsto\neg(A\equiv B)\\
   \neg A,B\mapsto\neg(A\equiv B)&\neg A,\neg B\mapsto A\equiv B\end{array}}
  {All four Boolean responses; repeated query labels are different prefix
   occurrences. Binary premise retained in full scenario.
   \ref{annex:p18.es.search-procedures}.}


\begin{hbtheorem}{thm:p18-es-search-procedures}{proved=annex:p18.es.search-procedures}{All eight original procedures and their underlying scenarios}
Every response branch of source Figures 5--12 has the sheet's stated principal
answer under its premises and Boolean or first-order truth clauses. The three
compound short-circuit procedures have three paths each, equivalence has four,
and negation has two. Inserting the indicated binary question gives the actual
scenario warrants for Figures 9--12. The procedures do not claim that merely
accepting an unverified response makes their endpoint true.
\end{hbtheorem}

\subsection{Related methods and exact boundaries}

The final source section reports connections to Urba\'nski's Synthetic Tableaux
Method and separate reducibility results, and mentions applications to
explanation by specification, abnormic hypotheses and reconstructions of
scientific discovery \cite[pp.~40--43]{Wisniewski2001QuestionsInferences}.
Those reports do not supply the referred tableaux rules, completeness proofs
or application studies here. A search-rule procedure is not by itself such
a proof system. The source's classical omniscience and inconsistent-knowledge
limitations remain: semantic non-entailment is not a model of everything a
human has yet failed to discover, and classical evocation excludes an
inconsistent premise set. The cited adaptive alternatives require their own
actual constructions; none is silently supplied by this scenario treatment.
